\usepackage[linesnumbered,algoruled,boxed,lined]{algorithm2e}
\usepackage{complexity} 
\usepackage{amssymb,amsmath,amsthm}
\usepackage{thm-restate}
\usepackage{mathtools}
\usepackage[a4paper, total={6.5in, 8.5in}]{geometry}
\usepackage{tikz}
\usetikzlibrary{fit}
\usepackage{hyperref}
\usepackage{cleveref}
\usepackage{cite}

\newcommand\addvmargin[1]{
  \node[fit=(current bounding box),inner ysep=#1,inner xsep=0]{};
}

\hypersetup{
    breaklinks=true,
    colorlinks=true,
    citecolor=red!70!blue!60!black,
    linkcolor=red!70!blue!90!black,
}

\usepackage[framemethod=tikz]{mdframed}
\mdfsetup{frametitlealignment=\centering}
\usepackage{enumitem}
\setlist[itemize]{noitemsep}
\setlist[itemize]{nosep}
\setlist[enumerate]{noitemsep}
\usepackage{vmargin}
\setmarginsrb{1in}{1in}{1in}{1in}{0pt}{0pt}{0pt}{6mm}

\usepackage{subcaption}
\usepackage{color}
\newcommand{\prelim}[1]{\textcolor{cyan}{#1}}
\newcommand{\tocheck}[1]{\textcolor{cyan}{#1}}
\newcommand{\todo}[1]{\textcolor{red}{\textbf{#1}}}
\newcommand{\blue}[1]{{\color{blue}{#1}}}
\newcommand{\red}[1]{{\color{red}[#1]}}
\newcommand{\rs}[1]{\textcolor{violet}{RS- #1}}
\newcommand{\eg}[1]{\textcolor{purple}{EG- #1 }}
\newcommand{\skb}[1]{\textcolor{pink!50!black}{SKB- #1}}
\newcommand{\dm}[1]{\textcolor{green}{DM- #1}}

\newcommand{\grid}{\boxplus}

\newcommand{\al}{\alpha}
\newcommand{\bt}{\beta}
\newcommand{\bd}{\partial}
\newcommand{\inter}{\mathrm{int}}

\newcommand{\calA}{\mathcal{A}}
\newcommand{\calB}{\mathcal{B}}
\newcommand{\calC}{\mathcal{C}}
\newcommand{\calD}{\mathcal{D}}
\newcommand{\calF}{\mathcal{F}}
\newcommand{\calG}{\mathcal{G}}
\newcommand{\calH}{\mathcal{H}}
\newcommand{\calO}{\mathcal{O}}
\newcommand{\calP}{\mathcal{P}}
\newcommand{\calR}{\mathcal{R}}
\newcommand{\calS}{\mathcal{S}}
\newcommand{\calT}{\mathcal{T}}
\newcommand{\calX}{\mathcal{X}}
\newcommand{\calW}{\mathcal{W}}
\newcommand{\calZ}{\mathcal{Z}}

\newcommand{\true}{\texttt{True}}
\newcommand{\false}{\texttt{False}}

\newcommand{\complete}{\textsf{complete}}
\newcommand{\ETH}{\textsf{ETH}}
\newcommand{\hard}{\textsf{hard}}
\newcommand{\NPpoly}{\textsf{NP/poly}}
\newcommand{\coNPpoly}{\textsf{coNP/poly}}
\newcommand{\paraNP}{\textsf{para-NP}}
\newcommand{\tw}{\mathrm{tw}}
\newcommand{\dist}{\mathrm{dist}}
\newcommand{\Skel}{\mathrm{Skel}}
\newcommand{\orig}{\mathrm{orig}}

\newcommand{\yes}{\textsc{Yes}}
\newcommand{\no}{\textsc{No}}
\newcommand{\out}{\texttt{out}}
\newcommand{\inp}{\texttt{in}}

\theoremstyle{plain}
\newtheorem{theorem}{Theorem}[section]
\newtheorem{conjecture}[theorem]{Conjecture}
\newtheorem{lemma}[theorem]{Lemma}
\newtheorem{corollary}[theorem]{Corollary}
\newtheorem{observation}[theorem]{Observation}
\newtheorem{claim}[theorem]{Claim}
\newtheorem{identification rule}{Identification Rule}
\newtheorem{invariants}{Invariants}
\newtheorem*{invariants*}{Invariants}


\theoremstyle{definition}
\newtheorem{definition}[theorem]{Definition}

\theoremstyle{remark}
\newtheorem{remark}[theorem]{Remark}

% ################ NEW ENVIROMENT STARTS ##########################
\newcommand{\defparproblem}[4]{
  \vspace{1mm}
\noindent\fbox{
  \begin{minipage}{0.96\textwidth}
  \begin{tabular*}{\textwidth}{@{\extracolsep{\fill}}lr} #1  & 
\\ \end{tabular*}
  {\bf{Input:}} #2  \\
  {\bf{Question:}} #4
  \end{minipage}
  }
  \vspace{1mm}
}

\newcommand{\defparproblemoutput}[4]{
  \vspace{1mm}
\noindent\fbox{
  \begin{minipage}{0.96\textwidth}
  \begin{tabular*}{\textwidth}{@{\extracolsep{\fill}}lr} #1  & {\bf{Parameter:}} #3
\\ \end{tabular*}
  {\bf{Input:}} #2  \\
  {\bf{Output:}} #4
  \end{minipage}
  }
  \vspace{1mm}
}

\newcommand{\defproblem}[3]{
  \vspace{1mm}
\noindent\fbox{
  \begin{minipage}{0.96\textwidth}
  \begin{tabular*}{\textwidth}{@{\extracolsep{\fill}}lr} #1 \\ \end{tabular*}
  {\bf{Input:}} #2  \\
  {\bf{Question:}} #3
  \end{minipage}
  }
  \vspace{1mm}
}

\newcommand{\defproblemout}[3]{
  \vspace{1mm}
\noindent\fbox{
  \begin{minipage}{0.96\textwidth}
  \begin{tabular*}{\textwidth}{@{\extracolsep{\fill}}lr} #1 \\ \end{tabular*}
  {\bf{Input:}} #2  \\
  {\bf{Output:}} #3
  \end{minipage}
  }
  \vspace{1mm}
}

\newcommand{\parnamedefn}[4]{
\begin{tabbing}
\name{#1} \\
\emph{Input:} \hspace{1cm} \= \parbox[t]{10cm}{#2} \\
\emph{Parameter:}            \> \parbox[t]{10cm}{#3} \\
\emph{Question:}             \> \parbox[t]{10cm}{#4} \\
\end{tabbing}
\vspace{-0.5cm}
}

\newcommand{\decnamedefn}[3]{
  \begin{tabbing} #1\\
    \emph{Input:} \hspace{1.2cm} \= \parbox[t]{12cm}{#2} \\
    \emph{Question:}             \> \parbox[t]{12cm}{#3} \\
  \end{tabbing}
}

\newcommand{\defparproblemnoname}[3]{
  \vspace{1mm}
\noindent\fbox{
  \begin{minipage}{0.96\textwidth}
  \begin{tabular*}{\textwidth}{@{\extracolsep{\fill}}lr} {\bf{Input:}} #1  & 
\\ \end{tabular*}
  {\bf{Question:}} #3
  \end{minipage}
  }
  \vspace{1mm}
}

\newenvironment{claimproof}{\begin{proof}\renewcommand{\qedsymbol}{\claimqed}}{\end{proof}\renewcommand{\qedsymbol}{\plainqed}}
\let\plainqed\qedsymbol
\newcommand{\claimqed}{$\lrcorner$}

% ################ NEW ENVIROMENT ENDS ##########################

\newcommand{\dsnfull}{{\sc Planar Directed Steiner Network}\xspace}
\newcommand{\dsngen}{{\sc Directed Steiner Network}\xspace}
\newcommand{\dsn}{{\sc Planar DSN}\xspace}
\newcommand{\scssfull}{{\sc Planar Strongly Connected Steiner Network}\xspace}
\newcommand{\scss}{{\sc Planar SCSS}\xspace}
\newcommand{\Dsnfull}{{\sc Planar $\calD$-Steiner Network}\xspace}
\newcommand{\Dsn}{{\sc Planar $\calD$-SN}\xspace}
\newcommand{\Dsngenfull}{{\sc $\calD$-Steiner Network}\xspace}
\newcommand{\Dsngen}{{\sc $\calD$-SN}\xspace}
\newcommand{\ADsnfull}{{\sc Acyclic Planar $\calD$-Steiner Network}\xspace}
\newcommand{\ADsn}{{\sc AP-$\calD$-SN}\xspace}
\newcommand{\minimal}{reachability-minimal}
\newcommand{\Minimal}{Reachability-minimal}

\newcommand{\tough}{tough}
\newcommand{\toughpair}{tough-pair}
\newcommand{\source}{source}
\newcommand{\head}{\text{head}}
\newcommand{\tail}{\text{tail}}
\newcommand{\type}[2]{\texttt{type}^{#2}(#1)}
\newcommand{\ind}{\ominus}
\newcommand{\outd}{\oplus}
\newcommand{\noth}{\oslash}

\newcommand{\oldC}{\mathcal{A}}

\newcommand{\minimalbiclique}{minimal-biclique}
\newcommand{\strictsc}{strict-strongly-connected}

\newcommand{\aux}[1]{\textnormal{\texttt{Aux}}_{#1}}
\newcommand{\col}{\texttt{col}}
\newcommand{\g}{g}

\newcommand{\contr}{\textnormal{\texttt{contr}}}
\newcommand{\Dcontr}{D_{\contr}}

\newcommand{\bip}{\texttt{bip}}

\newcommand{\sss}{\texttt{s}}
\newcommand{\ttt}{\texttt{t}}
\newcommand{\cleanedbiclique}{cleaned-biclique}

% ################ TIKZ ##########################

\usetikzlibrary{arrows}
\usetikzlibrary{arrows.meta}
\usetikzlibrary{decorations.pathmorphing}
\usetikzlibrary{decorations.markings}
\usetikzlibrary{calc}
\usetikzlibrary{positioning}

\definecolor{green}{RGB}{18, 122, 22}
\definecolor{red}{RGB}{219, 7, 7}
\definecolor{blue}{RGB}{0,0,204}
\tikzset{
  circ/.style = {circle,draw,fill,inner sep=1.3pt},
  mcirc/.style = {circle,draw,fill,inner sep=1pt},
  circR/.style = {circle,draw=red,fill=red,text=red,inner sep=1.3pt},
  circG/.style = {circle,draw=green,fill=green,text=green,inner sep=1.3pt},
  circB/.style = {circle,draw=blue,fill=blue,text=blue,inner sep=1.3pt},
  circb/.style = {circle,draw=blue,fill=blue,text=blue,inner sep=1.1pt},
  circr/.style = {circle,draw=red,fill=red,inner sep=1pt},
  scirc/.style = {circle,draw,fill,inner sep=.8pt},
  invisible/.style = {draw=none,inner sep=0pt,font=\tiny},
  nonedge/.style={decorate,decoration={snake,amplitude=.3mm,segment length=1mm},draw}
}

\newcommand{\executeiffilenewer}[3]{
\ifnum\pdfstrcmp{\pdffilemoddate{#1}}
{\pdffilemoddate{#2}}>0
{\immediate\write18{#3}}\fi
} 
\newcommand{\includesvg}[1]{
\executeiffilenewer{#1.svg}{#1.pdf}
{inkscape -z -D --file=#1.svg 
--export-pdf=#1.pdf --export-latex}
{\input{#1.pdf_tex}}}

\newcommand{\svg}[2]{\def\svgwidth{#1}\includesvg{#2}}

%%%%%%%%%%%%%Important macros%%%%%%%%%%%

\newcommand{\hardbiclique}[1]{$#1$-hard-biclique-pattern}
\newcommand{\hardmatching}[1]{$#1$-hard-matching-pattern}
\newcommand{\hardpattern}[1]{$#1$-hard-pattern}

\newcommand{\htt}{h(t)}
\newcommand{\corr}{\textnormal{\texttt{cor}}}
\newcommand{\cor}[1]{\corr(#1)}

%%%%%%%%%%%%%%%%%%%%%%

\newcommand{\aaa}{\texttt{a}}
\newcommand{\bbb}{\texttt{b}}